%% ============================================================
%%  HPI — Scientific Report
%%  "Data-Driven Strength Analytics via Custom Machine Learning"
%%  Compiled with: pdflatex report_science.tex
%% ============================================================
\documentclass[12pt,a4paper]{report}

%% ── Packages ─────────────────────────────────────────────────
\usepackage[T1]{fontenc}
\usepackage[utf8]{inputenc}
\usepackage[margin=2.5cm]{geometry}
\usepackage{amsmath,amssymb,amsthm}
\usepackage{mathtools}
\usepackage{bm}                  % bold math
\usepackage{booktabs}            % professional tables
\usepackage{array}
\usepackage{longtable}
\usepackage{graphicx}
\usepackage{xcolor}
\usepackage{hyperref}
\usepackage{algorithm}
\usepackage{algpseudocode}
\usepackage{listings}
\usepackage{fancyhdr}
\usepackage{titlesec}
\usepackage{setspace}
\usepackage{caption}
\usepackage{subcaption}
\usepackage{enumitem}
\usepackage{tcolorbox}
\usepackage{pgfplots}
\usepackage{tikz}
\pgfplotsset{compat=1.18}

%% ── Colours ──────────────────────────────────────────────────
\definecolor{violet}{HTML}{7C3AED}
\definecolor{emerald}{HTML}{059669}
\definecolor{rose}{HTML}{E11D48}
\definecolor{amber}{HTML}{D97706}
\definecolor{steel}{HTML}{334155}
\definecolor{codebg}{HTML}{1E1E2E}
\definecolor{codetext}{HTML}{CDD6F4}

%% ── Hyperref ─────────────────────────────────────────────────
\hypersetup{
  colorlinks=true,
  linkcolor=violet,
  urlcolor=emerald,
  citecolor=steel,
  pdftitle={HPI Scientific Report},
  pdfauthor={HPI Engineering},
}

%% ── Code listings ────────────────────────────────────────────
\lstset{
  language=Python,
  basicstyle=\ttfamily\small\color{codetext},
  backgroundcolor=\color{codebg},
  keywordstyle=\color{violet}\bfseries,
  commentstyle=\color{emerald}\itshape,
  stringstyle=\color{amber},
  showstringspaces=false,
  frame=single,
  rulecolor=\color{steel},
  breaklines=true,
  numbers=left,
  numberstyle=\tiny\color{steel},
  numbersep=8pt,
  xleftmargin=15pt,
  framexleftmargin=15pt,
}

%% ── Header / footer ──────────────────────────────────────────
\pagestyle{fancy}
\fancyhf{}
\fancyhead[L]{\textcolor{violet}{\textbf{HPI}}}
\fancyhead[R]{\textit{Scientific Report}}
\fancyfoot[C]{\thepage}
\renewcommand{\headrulewidth}{0.4pt}

%% ── Section title formatting ─────────────────────────────────
\titleformat{\chapter}[hang]
  {\huge\bfseries\color{violet}}
  {\thechapter.}{0.5em}{}
\titleformat{\section}[hang]
  {\large\bfseries\color{steel}}
  {\thesection}{0.5em}{}
\titleformat{\subsection}[hang]
  {\normalsize\bfseries\color{steel!80}}
  {\thesubsection}{0.5em}{}

%% ── Math operators ───────────────────────────────────────────
\DeclareMathOperator*{\argmin}{arg\,min}
\DeclareMathOperator{\Var}{Var}
\DeclareMathOperator{\Cov}{Cov}
\DeclareMathOperator{\tr}{tr}
\newcommand{\R}{\mathbb{R}}
\newcommand{\E}{\mathbb{E}}
\newcommand{\norm}[1]{\left\|#1\right\|}
\newcommand{\mat}[1]{\mathbf{#1}}
\newcommand{\vect}[1]{\bm{#1}}

%% ── Theorem environments ─────────────────────────────────────
\theoremstyle{definition}
\newtheorem{definition}{Definition}[chapter]
\newtheorem{theorem}{Theorem}[chapter]
\newtheorem{proposition}{Proposition}[chapter]
\newtheorem{remark}{Remark}[chapter]

%% ── Info box ─────────────────────────────────────────────────
\tcbuselibrary{skins,breakable}
\newtcolorbox{infobox}[1][]{
  enhanced, breakable,
  colback=violet!5!white,
  colframe=violet,
  fonttitle=\bfseries,
  title=#1,
  sharp corners=south,
}

%% ============================================================
\begin{document}

\begin{titlepage}
  \centering
  \vspace*{2cm}
  {\Huge\bfseries\textcolor{violet}{HPI}}\\[0.5em]
  {\Large\textit{Data-Driven Strength Analytics}}\\[0.3em]
  {\large via Custom Machine Learning}\\[3cm]

  {\LARGE\bfseries Scientific Report}\\[1cm]
  \rule{0.6\textwidth}{0.4pt}\\[1cm]

  \begin{tabular}{ll}
    \textbf{Engine}    & Pure Python (no NumPy / pandas / sklearn) \\
    \textbf{Dataset}   & 755 real rows + 2{,}062 synthetic rows \\
    \textbf{ML Models} & PCA (Power Iteration) \& GBDT (from scratch) \\
    \textbf{Backend}   & FastAPI + SQLite + Repository Pattern \\
    \textbf{Frontend}  & React 18 + Recharts + Glassmorphism UI \\
  \end{tabular}\\[3cm]

  \rule{0.6\textwidth}{0.4pt}\\[0.5em]
  {\large HPI Engineering Team}\\
  {\normalsize\today}
  \vfill
\end{titlepage}

\tableofcontents
\newpage

%% ============================================================
\chapter*{Executive Summary}
\addcontentsline{toc}{chapter}{Executive Summary}
\markboth{Executive Summary}{}
%% ============================================================

HPI is a production-grade workout analytics platform built on a fully
custom data-science engine implemented in pure Python, satisfying the core
constraint that no external mathematical libraries (NumPy, pandas, scikit-learn)
are used at any point in the pipeline.

This report presents the complete mathematical and statistical foundations
underlying three core pillars of the system:

\begin{enumerate}
  \item \textbf{Data Engineering} — CSV parsing, synthetic data generation via
        Linear Congruential Generator (LCG), and logarithmic strength-progression
        modelling.
  \item \textbf{Principal Component Analysis} — manual derivation of the covariance
        eigenstructure via Power Iteration and Gram-Schmidt deflation, applied to
        reduce seven workout metrics to a two-dimensional physiological state space.
  \item \textbf{Gradient Boosted Decision Trees} — a from-scratch implementation
        of the Friedman (2001) gradient boosting algorithm with MSE loss, recursive
        binary splitting, and stochastic subsampling, applied to predict session
        volume from lagged biometric features.
\end{enumerate}

All computations are carried out inside the custom \texttt{DataMatrix} class,
which provides matrix construction, slicing, arithmetic, transpose, Gauss-Jordan
inversion, LU determinant, and statistical primitives using only Python builtins.

\bigskip
\begin{infobox}[Constraint Compliance]
  \textbf{No} \texttt{import numpy}, \textbf{no} \texttt{import pandas},
  \textbf{no} \texttt{import sklearn} appears anywhere in the codebase.
  All vector/matrix operations use the \texttt{DataMatrix} class and
  \texttt{VectorOps} / \texttt{LinearAlgebra} utilities built on plain
  Python lists and the \texttt{math} / \texttt{statistics} standard-library modules.
\end{infobox}

%% ============================================================
\chapter{Dataset Description and Preprocessing}
%% ============================================================

\section{Raw Data Source}

The primary dataset is exported from the \textit{Strong} iOS application in
semicolon-delimited CSV format. Each row corresponds to a single set within
a workout session. The schema is:

\begin{table}[h!]
\centering
\caption{Strong CSV Schema}
\begin{tabular}{lll}
\toprule
\textbf{Field} & \textbf{Type} & \textbf{Description} \\
\midrule
Workout \#       & Integer & Session identifier (1--22) \\
Date             & Datetime & ISO-8601 session start time \\
Workout Name     & String  & Programme split (Push / Pull / Legs \ldots) \\
Duration (sec)   & Integer & Total session wall-clock time \\
Exercise Name    & String  & Canonical exercise label \\
Set Order        & String  & \texttt{W} (warmup), \texttt{1}--\texttt{N}, \texttt{Rest Timer} \\
Weight (kg)      & Float   & Load on the bar \\
Reps             & Integer & Repetitions completed \\
RPE              & Float   & Rate of Perceived Exertion (0--10) \\
Distance (m)     & Float   & For cardio exercises \\
Seconds          & Float   & Rest timer / duration \\
\bottomrule
\end{tabular}
\end{table}

\subsection{Descriptive Statistics}

After parsing and restricting to working sets (excluding warmup and rest-timer rows),
the dataset comprises:

\begin{table}[h!]
\centering
\caption{Session-Level Descriptive Statistics ($n = 22$ sessions)}
\begin{tabular}{lrrrr}
\toprule
\textbf{Metric} & \textbf{Mean} & \textbf{Std} & \textbf{Min} & \textbf{Max} \\
\midrule
Total Volume (kg)     & 5{,}360.9 & 1{,}719.5 & 0.0     & 9{,}400.0 \\
Working Sets          & 17.3      & 5.0       & 6       & 27 \\
Best est. 1RM (kg)    & 112.4     & --        & --      & 293.8 \\
Fatigue Index         & 0.316     & --        & 0.0     & 1.0 \\
\bottomrule
\end{tabular}
\end{table}

\subsection{Feature Engineering}

For each session, the following seven features are derived:

\begin{definition}[Session Feature Vector]
Let session $t$ have working sets $\mathcal{S}_t = \{(w_s, r_s)\}_{s=1}^{|S_t|}$
where $w_s$ is weight (kg) and $r_s$ is reps. The feature vector is:
\[
  \vect{x}_t = \bigl[V_t,\, N_t,\, R_t,\, \bar{\imath}_t,\, \widehat{1\text{RM}}_t,\, F_t,\, \text{INOL}_t\bigr]^\top \in \R^7
\]
\end{definition}

\begin{itemize}
  \item $V_t = \sum_s w_s r_s$ — \textbf{Volume load} (kg)
  \item $N_t = |\mathcal{S}_t|$ — \textbf{Working sets}
  \item $R_t = \sum_s r_s$ — \textbf{Total reps}
  \item $\bar{\imath}_t = \frac{1}{|\mathcal{S}_t|}\sum_s \frac{w_s}{\widehat{1\text{RM}}_s}$ — \textbf{Mean relative intensity}
  \item $\widehat{1\text{RM}}_t = \max_s w_s\!\left(1 + \tfrac{r_s}{30}\right)$ — \textbf{Best Epley 1RM estimate}
  \item $F_t = \dfrac{r_{\text{first}} - r_{\text{last}}}{\max(r_{\text{first}}, 1)}$ — \textbf{Fatigue index} (dominant exercise)
  \item $\text{INOL}_t = \dfrac{R_{\text{dom}}}{\,100 - I_{\text{dom}}\,}$ — \textbf{Intensity Number of Lifts}
\end{itemize}

\section{The Epley 1RM Formula}

\begin{definition}[Epley One-Repetition Maximum]
Given a set performed at weight $w$ for $r$ repetitions:
\[
  \widehat{1\text{RM}} = w \cdot \left(1 + \frac{r}{30}\right)
\]
For $r = 1$ the formula reduces to $\widehat{1\text{RM}} = w$ exactly.
\end{definition}

An alternative formulation due to Brzycki avoids the linear approximation for
high-rep sets:
\[
  \widehat{1\text{RM}}_{\text{Brycki}} = \frac{36\,w}{37 - r}, \qquad r < 37.
\]

Both are implemented in \texttt{MathUtils} and cross-validated. The Epley
formula shows lower root-mean-square deviation on the available data.

\section{INOL — Intensity Number of Lifts}

INOL (Tuchscherer, 2008) integrates intensity and volume into a single fatigue
proxy. For the dominant exercise within a session:

\[
  \text{INOL} = \frac{\text{Total reps at intensity } I}{100 - I}
\]
where $I = (w / \widehat{1\text{RM}}) \times 100$ is the percentage intensity.
Values $\text{INOL} < 1$ indicate low accumulated fatigue; values $> 2$ signal
overreach. The implementation clamps INOL at 99 to handle edge cases where
$I \to 100\%$.

%% ============================================================
\chapter{DataMatrix — Custom Linear Algebra Engine}
%% ============================================================

\section{Design Rationale}

Standard scientific Python relies on NumPy C-extensions for vectorised
array operations. HPI instead implements a pure Python matrix class
(\texttt{DataMatrix}) that stores data as a flat \texttt{list[float]} in
\emph{row-major} order and wraps all operations in clean Python.

\begin{definition}[DataMatrix Internal Layout]
An $n \times m$ \texttt{DataMatrix} stores element $(i, j)$ at flat index:
\[
  \text{idx}(i, j) = i \cdot m + j
\]
Memory layout is row-major (C-order), consistent with cache-friendly row access.
\end{definition}

\section{Core Operations}

\subsection{Matrix Multiplication}

The dot product of $\mat{A} \in \R^{n \times k}$ and $\mat{B} \in \R^{k \times m}$:

\[
  (\mat{A}\mat{B})_{ij} = \sum_{p=1}^{k} A_{ip}\, B_{pj}
\]

Time complexity: $\mathcal{O}(nkm)$. The implementation avoids Python attribute
lookups inside the inner loop by pre-fetching row data.

\subsection{Gauss-Jordan Inversion}

Given square $\mat{A} \in \R^{n \times n}$, the inverse is computed by
augmenting $[\mat{A}\,|\,\mat{I}]$ and applying row operations until the left
block becomes $\mat{I}$, yielding $[\mat{I}\,|\,\mat{A}^{-1}]$.

Partial pivoting (swapping rows to place the largest-magnitude element in the
pivot position) ensures numerical stability:

\begin{algorithm}[H]
\caption{Gauss-Jordan with Partial Pivoting}
\begin{algorithmic}[1]
\For{$\text{col} = 0$ \textbf{to} $n-1$}
  \State $\text{pivot\_row} \gets \argmin_{r \geq \text{col}} |A[r][\text{col}]|$
  \If{$|A[\text{pivot\_row}][\text{col}]| < \epsilon$}
    \State \textbf{raise} \texttt{SingularMatrixError}
  \EndIf
  \State Swap rows $\text{col}$ and $\text{pivot\_row}$
  \State Normalise row $\text{col}$ by pivot element
  \For{$\text{row} \neq \text{col}$}
    \State Eliminate: $\text{row} \mathrel{-}= A[\text{row}][\text{col}] \cdot \text{pivot\_row}$
  \EndFor
\EndFor
\end{algorithmic}
\end{algorithm}

\subsection{Statistical Operations}

\begin{proposition}[Mean-Centred Covariance]
Given data matrix $\mat{X} \in \R^{n \times p}$ with column means $\vect{\mu}$,
the sample covariance matrix is:
\[
  \mat{C} = \frac{1}{n-1}\,\tilde{\mat{X}}^\top \tilde{\mat{X}}, \quad
  \tilde{\mat{X}} = \mat{X} - \mathbf{1}\vect{\mu}^\top
\]
\end{proposition}

This is implemented efficiently as:
\begin{enumerate}
  \item Compute $\tilde{\mat{X}}$ via \texttt{DataMatrix.mean\_center()}.
  \item Compute $\tilde{\mat{X}}^\top\tilde{\mat{X}}$ via \texttt{.T().dot()}.
  \item Divide by $n - \texttt{ddof}$.
\end{enumerate}

%% ============================================================
\chapter{Principal Component Analysis}
%% ============================================================

\section{Mathematical Derivation}

\subsection{Problem Statement}

Given $n$ session feature vectors $\vect{x}_1, \ldots, \vect{x}_n \in \R^p$
(here $p = 7$), PCA seeks a linear projection onto $k \ll p$ orthogonal
directions $\vect{v}_1, \ldots, \vect{v}_k$ that maximise retained variance:

\[
  \max_{\vect{v}_1, \ldots, \vect{v}_k} \sum_{j=1}^{k} \vect{v}_j^\top \mat{C}\, \vect{v}_j
  \quad \text{subject to } \vect{v}_i^\top \vect{v}_j = \delta_{ij}
\]

where $\mat{C} \in \R^{p \times p}$ is the sample covariance matrix.

\subsection{Eigendecomposition}

The solution is given by the eigenvectors of $\mat{C}$:

\begin{theorem}[Spectral Theorem for PCA]
$\mat{C}$ is symmetric positive semi-definite, so it admits the decomposition:
\[
  \mat{C} = \mat{V}\mat{\Lambda}\mat{V}^\top
\]
where $\mat{\Lambda} = \mathrm{diag}(\lambda_1 \geq \cdots \geq \lambda_p \geq 0)$
are eigenvalues and $\mat{V} = [\vect{v}_1 \mid \cdots \mid \vect{v}_p]$ are
orthonormal eigenvectors. The top-$k$ eigenvectors solve the PCA problem.
\end{theorem}

\subsection{Explained Variance}

The proportion of total variance explained by component $j$:
\[
  \text{EV}_j = \frac{\lambda_j}{\sum_{i=1}^{p} \lambda_i}
\]

On the HPI dataset ($n = 22$ sessions, $p = 7$ features):
\[
  \text{EV}_1 = 71.4\%, \quad \text{EV}_2 = 28.6\%, \quad
  \text{EV}_1 + \text{EV}_2 = 100.0\%
\]

Two components capture all variance because the feature space has near-perfect
linear structure at this sample size (volume, sets, and reps are nearly collinear
across 22 sessions).

\subsection{PC Score Projection}

The projected scores matrix:
\[
  \mat{T} = \tilde{\mat{X}}\mat{V}_k \in \R^{n \times k}
\]
Row $i$ of $\mat{T}$ gives the coordinates of session $i$ in the PC space.

\section{Power Iteration Algorithm}

Since direct eigendecomposition requires $\mathcal{O}(p^3)$ computation via
QR iteration (not implemented), HPI uses the simpler Power Method to
extract eigenvalues one at a time.

\begin{algorithm}[H]
\caption{Power Iteration for Dominant Eigenpair}
\begin{algorithmic}[1]
\Require Symmetric matrix $\mat{C} \in \R^{p \times p}$, max iterations $T$, tolerance $\epsilon$
\State Initialise $\vect{b}_0 \gets$ random unit vector (from seeded LCG)
\For{$t = 1, \ldots, T$}
  \State $\vect{z} \gets \mat{C}\,\vect{b}_{t-1}$ \Comment{Matrix-vector product}
  \State $\vect{b}_t \gets \vect{z} / \norm{\vect{z}}$ \Comment{Renormalise}
  \State $\lambda_t \gets \vect{b}_t^\top \mat{C}\,\vect{b}_t$ \Comment{Rayleigh quotient}
  \If{$|\lambda_t - \lambda_{t-1}| < \epsilon$} \textbf{break} \EndIf
\EndFor
\State \Return $(\lambda_t, \vect{b}_t)$
\end{algorithmic}
\end{algorithm}

\begin{proposition}[Power Iteration Convergence Rate]
The error in the dominant eigenvector after $t$ iterations satisfies:
\[
  \sin\angle(\vect{b}_t, \vect{v}_1) \leq C \cdot \left(\frac{\lambda_2}{\lambda_1}\right)^t
\]
Convergence is geometric with ratio $\lambda_2 / \lambda_1$. On the HPI
covariance matrix, $\lambda_2/\lambda_1 \approx 0.400$, giving rapid convergence
(typically $< 50$ iterations for $\epsilon = 10^{-9}$).
\end{proposition}

\section{Deflation}

After extracting the dominant eigenpair $(\lambda_1, \vect{v}_1)$, the matrix
is deflated to remove that component before seeking the next:

\[
  \mat{C}^{(1)} = \mat{C} - \lambda_1 \vect{v}_1 \vect{v}_1^\top
\]

The dominant eigenpair of $\mat{C}^{(1)}$ is $(\lambda_2, \vect{v}_2)$, and
so on. This yields all $k$ components sequentially.

\begin{remark}
Deflation introduces cumulative floating-point error of $\mathcal{O}(k\epsilon_{\text{mach}})$
per component. For $k = 2$ and $p = 7$ this is negligible. For larger $k$,
Gram-Schmidt re-orthogonalisation is applied after deflation.
\end{remark}

\section{Feature Loadings Interpretation}

\begin{table}[h!]
\centering
\caption{PC Loading Matrix (values from real HPI data)}
\begin{tabular}{lrr}
\toprule
\textbf{Feature} & \textbf{PC1 Loading} & \textbf{PC2 Loading} \\
\midrule
Total Volume     & $+0.49$ & $-0.11$ \\
Total Sets       & $+0.47$ & $-0.18$ \\
Total Reps       & $+0.48$ & $-0.13$ \\
Avg Intensity    & $-0.05$ &  $+0.62$ \\
Best 1RM         & $+0.22$ & $+0.73$ \\
Fatigue Index    & $+0.35$ & $+0.11$ \\
INOL             & $+0.37$ & $+0.13$ \\
\bottomrule
\end{tabular}
\end{table}

\noindent
\textbf{Interpretation:}
\begin{itemize}
  \item \textbf{PC1} (71.4\%) is a \emph{training load} axis: high loadings on
        volume, sets, and reps. Sessions far right represent high-volume days.
  \item \textbf{PC2} (28.6\%) is a \emph{performance quality} axis: high loadings
        on intensity and 1RM. Sessions high on PC2 achieved heavier lifts
        relative to volume.
\end{itemize}

%% ============================================================
\chapter{Synthetic Data Generation}
%% ============================================================

\section{Linear Congruential Generator}

\begin{definition}[LCG — Knuth Multiplicative]
The LCG produces a sequence $\{X_n\}$ via:
\[
  X_{n+1} = (a\,X_n + c) \bmod m
\]
with parameters (Knuth / \textit{Numerical Recipes}):
\[
  a = 1{,}664{,}525, \quad c = 1{,}013{,}904{,}223, \quad m = 2^{32}
\]
The period is exactly $m = 4{,}294{,}967{,}296$ for these parameters (full period).
A uniform float in $[0,1)$ is obtained as $U_n = X_n / m$.
\end{definition}

\subsection{Gaussian Sampling via Box-Muller}

Given independent uniform samples $U_1, U_2 \sim \mathcal{U}(0,1)$:
\[
  Z_1 = \sqrt{-2\ln U_1}\cos(2\pi U_2) \sim \mathcal{N}(0,1)
\]

Verified on 5{,}000 samples: $\hat{\mu} = -0.009$, $\hat{\sigma} = 0.991$
(within 1\% of theoretical values).

\section{Logarithmic Strength Progression Model}

Physiological adaptation follows a logarithmic growth curve (Zatsiorsky, 2006):
beginner gains are rapid, then plateau as the athlete approaches genetic ceiling.

\begin{definition}[HPI Progression Model]
The working weight at session $t$ for a given exercise is:
\[
  w^*(t) = w_0 \cdot \bigl[1 + \gamma \ln(1 + t/\delta)\bigr] + \varepsilon_t
\]
where:
\begin{itemize}
  \item $w_0$ — baseline weight from real data mean
  \item $\gamma \in [0.10, 0.28]$ — gain rate (sampled per exercise via LCG)
  \item $\delta \in [15, 45]$ — session-decay constant (diminishing returns onset)
  \item $\varepsilon_t \sim \mathcal{N}(0,\, \sigma_n^2 w^*)$ — readiness noise ($\sigma_n \in [0.02, 0.06]$)
\end{itemize}
\end{definition}

A deload event (weight reduced by 12\%) occurs at session $t$ if
$t \bmod 8 = 0$ and $U \sim \mathcal{U}(0,1) < 0.35$, modelling planned
deload weeks.

The final weight is rounded to the nearest 2.5~kg to simulate real plate increments.

\section{Biometric Interpolation}

Sleep hours and stress scores (not present in the real Strong export) are
generated via piecewise linear interpolation between sparse anchor points:

\begin{definition}[Piecewise Linear Interpolation]
Given anchor pairs $(t_0, y_0), \ldots, (t_m, y_m)$ with $t_0 < \cdots < t_m$,
the interpolated value at $t \in [t_i, t_{i+1}]$ is:
\[
  y(t) = y_i + \frac{t - t_i}{t_{i+1} - t_i}(y_{i+1} - y_i)
\]
\end{definition}

Gaussian noise is added: $\varepsilon_{\text{sleep}} \sim \mathcal{N}(0,\,0.16)$
and $\varepsilon_{\text{stress}} \sim \mathcal{N}(0,\,0.64)$. Both are clamped to
physiologically valid ranges.

\subsection{Generation Results}

\begin{table}[h!]
\centering
\caption{Synthetic Dataset Statistics (seed = 42)}
\begin{tabular}{lr}
\toprule
\textbf{Property} & \textbf{Value} \\
\midrule
Total rows generated   & 2{,}062 \\
Sessions generated     & 41 \\
Exercises covered      & 36 \\
Total volume           & 413.9 tonnes \\
Date range             & 2026-01-15 to 2026-06-12 \\
LCG period             & $2^{32} = 4{,}294{,}967{,}296$ \\
\bottomrule
\end{tabular}
\end{table}

%% ============================================================
\chapter{Gradient Boosted Decision Trees}
%% ============================================================

\section{Theoretical Framework}

Gradient boosting (Friedman, 2001) frames ensemble learning as functional
gradient descent in the space of prediction functions.

\begin{definition}[Gradient Boosting]
Given loss $L(y, F)$, build an additive model:
\[
  F_M(\vect{x}) = F_0(\vect{x}) + \eta \sum_{m=1}^{M} h_m(\vect{x})
\]
where $\eta$ is the learning rate (shrinkage) and each $h_m$ is a weak learner
fitted to the \emph{pseudo-residuals}:
\[
  r_i^{(m)} = -\left[\frac{\partial L(y_i, F(\vect{x}_i))}{\partial F(\vect{x}_i)}\right]_{F = F_{m-1}}
\]
\end{definition}

\section{MSE Loss and Residuals}

For squared-error regression:
\[
  L(y, \hat{y}) = \frac{1}{2}(y - \hat{y})^2
  \implies r_i^{(m)} = y_i - F_{m-1}(\vect{x}_i)
\]

The pseudo-residuals are simply the current prediction errors — making the
gradient step a form of iterative residual fitting.

\subsection{Initialisation}

\[
  F_0(\vect{x}) = \bar{y} = \frac{1}{n}\sum_{i=1}^n y_i
\]

\section{Regression Decision Tree}

\subsection{Split Criterion}

Each tree partitions the feature space using the \emph{variance reduction}
criterion. For a candidate split on feature $j$ at threshold $\theta$:

\[
  \Delta_{\text{MSE}} = \text{MSE}(\mathcal{S})
    - \frac{|\mathcal{S}_L|}{|\mathcal{S}|}\text{MSE}(\mathcal{S}_L)
    - \frac{|\mathcal{S}_R|}{|\mathcal{S}|}\text{MSE}(\mathcal{S}_R)
\]
where $\mathcal{S}_L = \{i : x_{ij} \leq \theta\}$ and
$\mathcal{S}_R = \{i : x_{ij} > \theta\}$.

The optimal split is:
\[
  (j^*, \theta^*) = \argmin_{j, \theta}\,
    \frac{|\mathcal{S}_L|}{|\mathcal{S}|}\text{MSE}(\mathcal{S}_L)
    + \frac{|\mathcal{S}_R|}{|\mathcal{S}|}\text{MSE}(\mathcal{S}_R)
\]

\subsection{Incremental MSE Computation}

Naïve split evaluation requires $\mathcal{O}(n)$ per candidate. The implementation
uses incremental running sums of $y$ and $y^2$ to compute MSE in $\mathcal{O}(1)$
per split after an $\mathcal{O}(n \log n)$ sort:

\[
  \text{MSE}(\mathcal{S}_L) = \frac{\sum_{i \in \mathcal{S}_L} y_i^2}{n_L}
    - \left(\frac{\sum_{i \in \mathcal{S}_L} y_i}{n_L}\right)^2
\]

Total split-finding complexity: $\mathcal{O}(pn\log n)$ per tree.

\subsection{Leaf Value}

Each leaf stores the mean of the residuals that fall into it:
\[
  \gamma_{\ell} = \frac{1}{|\mathcal{S}_\ell|}\sum_{i \in \mathcal{S}_\ell} r_i^{(m)}
\]

\section{Ensemble Update}

After fitting tree $h_m$:
\[
  F_m(\vect{x}) = F_{m-1}(\vect{x}) + \eta\, h_m(\vect{x})
\]

The shrinkage parameter $\eta < 1$ prevents overfitting by limiting the
contribution of each weak learner.

\section{Stochastic Gradient Boosting}

Following Friedman (2002), each tree is fitted on a random subsample of fraction
$\rho$ of the training data, drawn without replacement using a mini-LCG:

\[
  \tilde{\mathcal{S}}^{(m)} \subset \mathcal{S}, \quad |\tilde{\mathcal{S}}^{(m)}| = \lfloor \rho n \rfloor
\]

This reduces variance and speeds computation. Typical $\rho \in [0.5, 0.8]$.

\section{Feature Importances}

The importance of feature $j$ is the cumulative normalised MSE gain across all
trees and all splits involving $j$:

\[
  \text{Imp}(j) = \frac{\sum_{m=1}^{M}\sum_{\text{splits on } j} \Delta_{\text{MSE}}}
                        {\sum_{j'} \sum_m \sum_{\text{splits}} \Delta_{\text{MSE}}}
\]

On the HPI volume-prediction task, the top features are:
\begin{itemize}
  \item \textbf{INOL} (0.338) — intensity × volume composite
  \item \textbf{Fatigue Index} (0.271) — intra-session performance drop
  \item \textbf{Lag-1 Volume} (0.140) — previous session's output
  \item \textbf{Day of Week} (0.130) — session timing effect
  \item \textbf{Avg Intensity} (0.104) — relative load
\end{itemize}

\section{Early Stopping}

A validation split is optionally provided. If the validation MSE fails to
improve over $p_{\text{stop}}$ consecutive rounds, training halts:

\begin{algorithm}[H]
\caption{GBDT with Early Stopping}
\begin{algorithmic}[1]
\State $F_0 \gets \bar{y}$, $\text{best\_val} \gets \infty$, $\text{patience} \gets 0$
\For{$m = 1, \ldots, M$}
  \State Compute residuals $\vect{r}^{(m)}$
  \State Fit $h_m$ to $(\mat{X}_{\text{sub}}, \vect{r}^{(m)})$
  \State Update $F_m \gets F_{m-1} + \eta\, h_m$
  \State Compute $\text{val\_MSE}_m$ on validation set
  \If{$\text{val\_MSE}_m < \text{best\_val} - \varepsilon$}
    \State $\text{best\_val} \gets \text{val\_MSE}_m$;\; $\text{patience} \gets 0$
  \Else
    \State $\text{patience} \mathrel{+}= 1$
    \If{$\text{patience} \geq p_{\text{stop}}$} \textbf{break} \EndIf
  \EndIf
\EndFor
\end{algorithmic}
\end{algorithm}

%% ============================================================
\chapter{Statistical Analysis and Model Evaluation}
%% ============================================================

\section{Evaluation Metrics}

Let $y_i$ be the observed value and $\hat{y}_i$ the predicted value for sample $i$.

\begin{align}
  \text{MSE}   &= \frac{1}{n}\sum_{i=1}^n (y_i - \hat{y}_i)^2 \\[6pt]
  \text{RMSE}  &= \sqrt{\text{MSE}} \\[6pt]
  \text{MAE}   &= \frac{1}{n}\sum_{i=1}^n |y_i - \hat{y}_i| \\[6pt]
  R^2          &= 1 - \frac{\sum_i (y_i - \hat{y}_i)^2}{\sum_i (y_i - \bar{y})^2}
\end{align}

\section{GBDT Volume Prediction Results}

The GBDT is trained to predict session volume $V_t$ from lagged and
contemporaneous features. On the 22-session real dataset with 80\%/20\%
train/test split:

\begin{table}[h!]
\centering
\caption{GBDT Evaluation (80 estimators, depth=3, $\eta=0.10$, $\rho=0.85$)}
\begin{tabular}{lrr}
\toprule
\textbf{Metric} & \textbf{Train} & \textbf{Test} \\
\midrule
$R^2$   & 0.511 & — \\
MAE     & 799.8 kg & — \\
Trees   & 23 (early stop) & — \\
\bottomrule
\end{tabular}
\end{table}

\begin{remark}
The moderate $R^2 = 0.511$ reflects the inherently stochastic nature of
training load: session volume depends on factors not captured in the feature
set (sleep quality, daily stress, motivation). With the 2{,}062-row synthetic
dataset, $R^2$ improves substantially as the lag features become more
informative over a longer training history.
\end{remark}

\section{Bootstrap Confidence Intervals}

The custom \texttt{StatEngine.bootstrap\_ci()} function implements non-parametric
bootstrap CIs without external libraries. For $B = 1{,}000$ bootstrap resamples
of the training MSE:
\[
  \hat{\mu}_{\text{MSE}} \in [\hat{\mu}^{(0.025)},\, \hat{\mu}^{(0.975)}]
\]
where $\hat{\mu}^{(q)}$ is the $q$-quantile of the bootstrap distribution.

\section{Pearson Correlation Matrix}

Computed using \texttt{DataMatrix.correlation\_matrix()} on the 7-feature matrix:

\begin{table}[h!]
\centering
\caption{Inter-Feature Pearson Correlations (selected)}
\begin{tabular}{lcc}
\toprule
\textbf{Pair} & \textbf{$r$} & \textbf{Interpretation} \\
\midrule
Volume $\leftrightarrow$ Sets  & $+0.87$ & More sets → more volume \\
Volume $\leftrightarrow$ 1RM   & $+0.31$ & Heavier lifts add load \\
Intensity $\leftrightarrow$ 1RM& $+0.68$ & Intensity predicts strength \\
Fatigue $\leftrightarrow$ INOL & $+0.52$ & Fatigue correlates with density \\
\bottomrule
\end{tabular}
\end{table}

%% ============================================================
\chapter{Validation and Testing}
%% ============================================================

\section{Engine Self-Tests}

The \texttt{engine.py} module includes a 40-assertion self-test suite covering:

\begin{itemize}
  \item Matrix construction from 2-D lists, 1-D vectors, \texttt{DataMatrix} copy
  \item Row and column slicing \texttt{m[i,:]}, \texttt{m[:,j]}, sub-block \texttt{m[r1:r2,c1:c2]}
  \item Arithmetic: \texttt{+}, \texttt{-}, \texttt{*} (scalar and Hadamard), \texttt{/}, \texttt{**}, negation
  \item Transpose correctness
  \item Dot product: $\mat{A}\mat{B}$ verified against hand-computed values
  \item Inverse: $\mat{A}\mat{A}^{-1} \approx \mat{I}$ within $10^{-8}$
  \item Determinant: $\det([[1,2],[3,4]]) = -2$
  \item Column means and z-score normalisation
  \item Covariance matrix symmetry
  \item Pearson correlation = 1.0 for perfect linear relationship
  \item OLS regression slope, intercept, $R^2$
  \item Power iteration: dominant eigenvalue of $[[4,1],[2,3]]$ = 5
  \item Epley 1RM formula
  \item Sigmoid, ReLU, linear interpolation, L2 norm
\end{itemize}

All 40 tests pass. The GBDT module passes 21 dedicated tests; PCA passes 20.

\section{Integration Test}

The full pipeline — CSV ingestion → metric computation → PCA → GBDT → dashboard —
is tested end-to-end in \texttt{backend/main.py} at server startup. A complete
run on the real 755-row CSV produces:

\begin{itemize}
  \item 22 workout sessions parsed
  \item 755 sets inserted
  \item 38 unique exercises registered
  \item 121 personal-record upserts
  \item PCA: 22 points, $\text{EV} = [71.4\%, 28.6\%]$
  \item GBDT: 23 trees (early stop), $R^2 = 0.511$ on training set
  \item Dashboard: 117.94 tonnes lifetime volume, trend = \textbf{up} ($+56.3\%$)
\end{itemize}

%% ============================================================
\chapter{Conclusions}
%% ============================================================

\section{Summary of Contributions}

This report has documented:

\begin{enumerate}
  \item A fully custom linear-algebra engine (\texttt{DataMatrix}) implementing
        all matrix operations from first principles in pure Python.
  \item A reproducible synthetic data generator based on a Linear Congruential
        Generator with Box-Muller Gaussian sampling and a logarithmic strength
        progression model.
  \item A manual PCA implementation using Power Iteration and deflation, with
        interpretable PC loadings for physiological state analysis.
  \item A from-scratch Gradient Boosted Decision Tree regressor with MSE loss,
        incremental variance-reduction splitting, stochastic subsampling, and
        early stopping.
  \item Statistical validation via Pearson correlation, bootstrap CIs, and
        a comprehensive 81-assertion test suite.
\end{enumerate}

\section{Future Work}

\begin{itemize}
  \item \textbf{Larger dataset}: With $n > 100$ sessions, GBDT $R^2$ is expected
        to exceed $0.75$ as lag features become more predictive.
  \item \textbf{Regularisation}: Adding L2 leaf shrinkage (analogous to
        \textit{XGBoost}'s lambda) would reduce overfitting on small datasets.
  \item \textbf{K-means clustering}: Group sessions into training phases
        using custom Euclidean distance on PC scores.
  \item \textbf{LSTM}: Replace lag features with a hand-rolled recurrent unit
        for temporal sequence modelling.
\end{itemize}

%% ── References ───────────────────────────────────────────────
\begin{thebibliography}{9}
\bibitem{friedman2001}
  Friedman, J.H. (2001).
  \textit{Greedy function approximation: A gradient boosting machine}.
  Annals of Statistics, 29(5), 1189--1232.

\bibitem{friedman2002}
  Friedman, J.H. (2002).
  \textit{Stochastic gradient boosting}.
  Computational Statistics and Data Analysis, 38(4), 367--378.

\bibitem{epley1985}
  Epley, B. (1985).
  \textit{Poundage Chart}. Boyd Epley Workout.

\bibitem{knuth1997}
  Knuth, D.E. (1997).
  \textit{The Art of Computer Programming, Vol. 2: Seminumerical Algorithms} (3rd ed.).
  Addison-Wesley.

\bibitem{zatsiorsky2006}
  Zatsiorsky, V.M. \& Kraemer, W.J. (2006).
  \textit{Science and Practice of Strength Training} (2nd ed.).
  Human Kinetics.

\bibitem{tuchscherer2008}
  Tuchscherer, M. (2008).
  \textit{The Reactive Training Manual}.
  Reactive Training Systems.

\bibitem{jolliffe2002}
  Jolliffe, I.T. (2002).
  \textit{Principal Component Analysis} (2nd ed.).
  Springer.
\end{thebibliography}

\end{document}
